% ============================================================
%  Three-Phase Phasor Extraction via DFT – Mathematical Notes
%  Author : Davide Bagnara
% ============================================================
\documentclass[11pt,a4paper]{article}

% --- Packages ---
\usepackage[T1]{fontenc}
\usepackage[utf8]{inputenc}
\usepackage[english]{babel}
\usepackage{lmodern}
\usepackage{amsmath, amssymb, amsthm}
\usepackage{mathtools}
\usepackage{bm}
\usepackage{geometry}
\usepackage{booktabs}
\usepackage{xcolor}
\usepackage{listings}
\usepackage{caption}
\usepackage{hyperref}
\usepackage{microtype}
\usepackage{parskip}

\geometry{top=2.5cm, bottom=2.5cm, left=2.8cm, right=2.8cm}

% --- Listings style for MATLAB ---
\definecolor{matlabgreen}{rgb}{0.13,0.55,0.13}
\definecolor{matlabblue} {rgb}{0.00,0.00,0.80}
\definecolor{matlabgray} {rgb}{0.50,0.50,0.50}

\lstdefinestyle{matlab}{
  language        = Matlab,
  basicstyle      = \ttfamily\small,
  keywordstyle    = \color{matlabblue}\bfseries,
  commentstyle    = \color{matlabgreen}\itshape,
  stringstyle     = \color{red},
  numbers         = left,
  numberstyle     = \tiny\color{matlabgray},
  numbersep       = 8pt,
  frame           = single,
  framerule       = 0.4pt,
  rulecolor       = \color{matlabgray},
  breaklines      = true,
  captionpos      = b,
  tabsize         = 4,
  showstringspaces= false,
}

% --- Math operators ---
\DeclareMathOperator{\Real}{Re}
\DeclareMathOperator{\Imag}{Im}
\newcommand{\jj}{\mathrm{j}}          % imaginary unit
\newcommand{\e}{\mathrm{e}}           % Euler's number
\newcommand{\DFT}[1]{\mathcal{F}\!\left\{#1\right\}}

% --- Theorem-like environments ---
\newtheorem{remark}{Remark}
\newtheorem{property}{Property}

% ============================================================
\begin{document}

\begin{center}
  {\LARGE\bfseries Three-Phase Phasor Extraction via DFT}\\[0.6em]
  {\large Mathematical Derivation and MATLAB Implementation}\\[0.4em]
  {\normalsize claude}\\[0.2em]
  {\small \today}
\end{center}

\vspace{0.5em}
\hrule
\vspace{1em}

% ============================================================
\section{Problem Statement}

Consider a balanced three-wire three-phase system whose line currents satisfy
Kirchhoff's current law at every instant:
\begin{equation}
  i_u(t) + i_v(t) + i_w(t) = 0 \qquad \forall\, t.
  \label{eq:kcl}
\end{equation}
In the fundamental case each current is a sinusoid of amplitude $\hat{I}$,
frequency $f_0$, and initial phase $\varphi$:
\begin{align}
  i_u(t) &= \hat{I}\cos(2\pi f_0 t + \varphi),              \label{eq:iu}\\
  i_v(t) &= \hat{I}\cos\!\left(2\pi f_0 t + \varphi - \tfrac{2\pi}{3}\right),\label{eq:iv}\\
  i_w(t) &= \hat{I}\cos\!\left(2\pi f_0 t + \varphi + \tfrac{2\pi}{3}\right).\label{eq:iw}
\end{align}
The constraint~\eqref{eq:kcl} is automatically verified because
$\cos\theta + \cos(\theta-2\pi/3) + \cos(\theta+2\pi/3) = 0$ for all $\theta$.

The goal is to represent each current as a \emph{phasor}
$\underline{I}_x \in \mathbb{C}$ in the polar plane, and to extract those
phasors numerically from sampled data using the Discrete Fourier Transform (DFT).

% ============================================================
\section{Phasor Representation}

A real sinusoid $x(t) = \hat{X}\cos(2\pi f_0 t + \phi)$ is identified by its
\emph{complex phasor}
\begin{equation}
  \underline{X} \;=\; \hat{X}\,\e^{\jj\phi} \;=\; \hat{X}\angle\phi,
  \label{eq:phasor_def}
\end{equation}
so that $x(t) = \Real\!\left[\underline{X}\,\e^{\jj 2\pi f_0 t}\right]$.

For the three-phase system above:
\begin{equation}
  \underline{I}_u = \hat{I}\e^{\jj\varphi},\qquad
  \underline{I}_v = \hat{I}\e^{\jj(\varphi - 2\pi/3)},\qquad
  \underline{I}_w = \hat{I}\e^{\jj(\varphi + 2\pi/3)}.
  \label{eq:phasors}
\end{equation}
The three phasors are displaced by $120°$ in the complex plane and satisfy
$\underline{I}_u + \underline{I}_v + \underline{I}_w = 0$.

% ============================================================
\section{Discrete Fourier Transform – Brief Review}

Let $x[n] = x(n/F_s)$, $n = 0,1,\ldots,N-1$, be a uniformly sampled sequence
with sampling frequency $F_s$ and $N$ total samples.  The DFT is defined as
\begin{equation}
  X[k] \;=\; \sum_{n=0}^{N-1} x[n]\,\e^{-\jj 2\pi k n / N},
  \qquad k = 0, 1, \ldots, N-1.
  \label{eq:dft}
\end{equation}
The DFT bin $k$ corresponds to the physical frequency
\begin{equation}
  f_k = \frac{k\,F_s}{N}.
  \label{eq:fk}
\end{equation}
For a pure cosine at $f_0$ sampled over an \emph{integer number of periods}
(i.e.\ $N = m\,F_s/f_0$ for some integer $m$), only two bins are non-zero:
$k_1 = m$ (positive frequency) and $k_{N-k_1}$ (negative frequency).  The
positive-frequency bin evaluates to
\begin{equation}
  X[k_1] = \frac{N}{2}\,\hat{X}\,\e^{\jj\phi},
  \label{eq:dft_cosine}
\end{equation}
so the phasor can be recovered as
\begin{equation}
  \boxed{%
  \underline{X} = \frac{2}{N}\,X[k_1]
               = \frac{2}{N}\sum_{n=0}^{N-1} x[n]\,\e^{-\jj 2\pi k_1 n/N}.}
  \label{eq:phasor_extract}
\end{equation}

\begin{remark}
  The factor $2/N$ accounts for the one-sided normalisation: the DFT at the
  positive-frequency bin captures half the total signal energy (the other half
  is in the conjugate bin at $N-k_1$).
\end{remark}

% ============================================================
\section{Derivation of the \texttt{phasor()} Function}
\label{sec:derivation}

Equation~\eqref{eq:phasor_extract} is the theoretical backbone of the
one-liner used in the MATLAB script:
\begin{equation}
  \texttt{phasor(x)} \;\equiv\;
  \frac{2}{N}\sum_{n=0}^{N-1} x[n]\,\e^{-\jj 2\pi k_1 n/N},
  \label{eq:phasor_matlab}
\end{equation}
where $k_1 = \texttt{round}(f_0 N / F_s)$ is the DFT bin index closest to the
fundamental frequency $f_0$.

\subsection*{Step-by-step derivation}

\noindent\textbf{Step 1 – Represent the sinusoid as a sum of complex
exponentials.}

Using Euler's formula $\cos\theta = \tfrac{1}{2}(\e^{\jj\theta}+\e^{-\jj\theta})$,
a single-tone signal of amplitude $\hat{X}$ and phase $\phi$ is
\begin{equation}
  x[n] = \hat{X}\cos\!\left(\frac{2\pi k_1 n}{N}+\phi\right)
       = \frac{\hat{X}}{2}\,\e^{\jj\phi}\,\e^{\jj 2\pi k_1 n/N}
       + \frac{\hat{X}}{2}\,\e^{-\jj\phi}\,\e^{-\jj 2\pi k_1 n/N}.
  \label{eq:euler}
\end{equation}

\noindent\textbf{Step 2 – Substitute into the DFT inner product at bin $k_1$.}

Multiplying $x[n]$ by the DFT kernel $\e^{-\jj 2\pi k_1 n/N}$ and summing:
\begin{align}
  \sum_{n=0}^{N-1} x[n]\,\e^{-\jj 2\pi k_1 n/N}
  &= \frac{\hat{X}}{2}\,\e^{\jj\phi}
     \underbrace{\sum_{n=0}^{N-1}\e^{\jj 2\pi k_1 n/N}\e^{-\jj 2\pi k_1 n/N}}_{= N}
  \notag\\
  &\quad + \frac{\hat{X}}{2}\,\e^{-\jj\phi}
     \underbrace{\sum_{n=0}^{N-1}\e^{-\jj 2\pi \cdot 2k_1 n/N}}_{= 0\text{ (orthogonality)}}.
  \label{eq:inner}
\end{align}
The second sum vanishes because $2k_1/N$ is not an integer multiple of $1/N$
(under the assumption $k_1 < N/2$), so the complex exponential completes full
cycles that cancel.  Hence
\begin{equation}
  \sum_{n=0}^{N-1} x[n]\,\e^{-\jj 2\pi k_1 n/N}
  = \frac{N}{2}\,\hat{X}\,\e^{\jj\phi}.
  \label{eq:inner_result}
\end{equation}

\noindent\textbf{Step 3 – Normalise to obtain the phasor.}

Multiplying both sides by $2/N$:
\begin{equation}
  \frac{2}{N}\sum_{n=0}^{N-1} x[n]\,\e^{-\jj 2\pi k_1 n/N}
  = \hat{X}\,\e^{\jj\phi} = \underline{X}. \qquad \square
  \label{eq:result}
\end{equation}
This is exactly the complex phasor defined in~\eqref{eq:phasor_def}, with
modulus $|\underline{X}| = \hat{X}$ (peak value) and argument
$\arg(\underline{X}) = \phi$.

\begin{remark}[Validity for non-pure-tone signals]
  If the signal contains multiple harmonics or noise, the DFT at bin $k_1$
  still isolates the fundamental component, because the DFT basis vectors are
  orthogonal.  Contributions from other frequencies land on different bins and
  do not corrupt $X[k_1]$.  This makes the method robust to distortion.
\end{remark}

% ============================================================
\section{MATLAB Implementation}

The derivation translates verbatim into the following anonymous function:

\begin{lstlisting}[style=matlab, caption={Phasor extraction via single-bin DFT (MATLAB).}]
% Number of samples in one fundamental period
N  = length(t);

% DFT bin index corresponding to f0
k1 = round(f0 * N / Fs);

% Anonymous function: single-bin DFT, normalised to peak phasor
phasor = @(x) (2/N) * sum( x .* exp(-1j * 2*pi * (0:N-1) * k1 / N) );

% Apply to each phase current
Iu = phasor(iu);   % Complex phasor of i_u(t)
Iv = phasor(iv);   % Complex phasor of i_v(t)
Iw = phasor(iw);   % Complex phasor of i_w(t)
\end{lstlisting}

\subsection*{Line-by-line explanation}

\begin{description}
  \item[\texttt{N = length(t)}]
    Total number of samples.  Since \texttt{t} spans exactly one period
    $T = 1/f_0$ with step $1/F_s$, we have $N = F_s/f_0$, ensuring integer
    DFT bin alignment.

  \item[\texttt{k1 = round(f0*N/Fs)}]
    Computes the fundamental bin index from~\eqref{eq:fk}:
    $k_1 = f_0 N / F_s$.  The \texttt{round} guards against floating-point
    rounding when $F_s/f_0$ is not exactly integer.

  \item[\texttt{exp(-1j*2*pi*(0:N-1)*k1/N)}]
    Evaluates the DFT kernel $\e^{-\jj 2\pi k_1 n/N}$ for all sample indices
    $n = 0,1,\ldots,N-1$ in a single vectorised operation.

  \item[\texttt{x .* (\ldots)}]
    Element-wise multiplication of the signal vector by the kernel vector,
    implementing the inner product $\sum_n x[n]\,\e^{-\jj 2\pi k_1 n/N}$.

  \item[\texttt{sum(\ldots)}]
    Accumulates the inner product over all $N$ samples.

  \item[\texttt{(2/N) * (\ldots)}]
    Applies the normalisation factor from~\eqref{eq:phasor_extract}, returning
    the complex phasor $\underline{X} = \hat{X}\e^{\jj\phi}$.
\end{description}

% ============================================================
\section{Verification}

For the balanced system~\eqref{eq:iu}--\eqref{eq:iw} with $\hat{I}=100\,\text{A}$
and $\varphi = 0$:
\begin{align*}
  \underline{I}_u &= 100\angle 0°,\\
  \underline{I}_v &= 100\angle{-120°},\\
  \underline{I}_w &= 100\angle{+120°}.
\end{align*}
These satisfy $|\underline{I}_u|=|\underline{I}_v|=|\underline{I}_w|=\hat{I}$
and $\underline{I}_u+\underline{I}_v+\underline{I}_w=0$, consistent with
the Kirchhoff constraint~\eqref{eq:kcl}.

\begin{property}[Symmetry check in MATLAB]
  The residual \texttt{max(abs(iu+iv+iw))} should be of order $10^{-12}$\,A
  (double-precision machine epsilon times $\hat{I}$), confirming that the
  constraint is satisfied to numerical precision.
\end{property}

% ============================================================
\section{Summary}

\begin{center}
\begin{tabular}{@{}lll@{}}
  \toprule
  Quantity & Symbol & MATLAB expression \\
  \midrule
  Peak amplitude   & $\hat{X}$              & \texttt{abs(phasor(x))} \\
  Phase angle      & $\phi$                 & \texttt{angle(phasor(x))} \\
  Real part        & $\hat{X}\cos\phi$      & \texttt{real(phasor(x))} \\
  Imaginary part   & $\hat{X}\sin\phi$      & \texttt{imag(phasor(x))} \\
  \bottomrule
\end{tabular}
\end{center}

The single-bin DFT approach is preferred over peak-detection or zero-crossing
methods because it is robust to harmonics, noise, and non-integer periods, and
because it generalises directly to unbalanced or distorted waveforms without
any modification to the code.

\end{document}
